% ============================================================================
% CHEATSHEET TEMPLATE
% ============================================================================
% Based on the ANA-G cheatsheet style
%
% Features:
% - 4-column landscape layout on A4
% - Colored boxes for main concepts and sub-concepts
% - Compact formatting optimized for information density
% - Mathematical typesetting support
%
% Usage:
% 1. Replace metadata (title, author, date)
% 2. Adjust number of columns if needed (see \Ncolumn)
% 3. Add your content using mainbox and subbox environments
% 4. Compile with pdflatex (run twice for proper layout)
% ============================================================================

\documentclass[a4paper, fontsize=7pt]{scrartcl}

% ============================================================================
% ENCODING & LANGUAGE
% ============================================================================
\usepackage[utf8]{inputenc}
% \usepackage[german]{babel}    % Uncomment for German hyphenation
% \usepackage[english]{babel}   % Uncomment for English hyphenation

% ============================================================================
% PAGE LAYOUT
% ============================================================================
% 4-column landscape layout with minimal margins
\usepackage[landscape, left=0.75cm, top=1cm, right=0.75cm, bottom=1.5cm, footskip=5pt]{geometry}
\usepackage{flowfram}
\ffvadjustfalse
\setlength{\columnsep}{1cm}

% Number of columns (adjust as needed: 3-4 recommended for A4)
\Ncolumn{4}

% Remove paragraph indentation
\setlength{\parindent}{0cm}

% ============================================================================
% COLORS & BOXES
% ============================================================================
\usepackage[most]{tcolorbox}

% Base style for all boxes
\tcbset{
  base/.style={
    boxrule=0.3mm,
    leftrule=1mm,
    left=1.75mm,
    arc=0mm,
    fonttitle=\bfseries,
    colbacktitle=black!10!white,
    coltitle=black,
    toptitle=0.75mm,
    bottomtitle=0.25mm,
    title={#1}
  }
}

% Main concept box (green accent)
% Use for: definitions, theorems, important formulas
\definecolor{maincolor}{rgb}{0.004, 0.5, 0.35}  % Nice green
\newtcolorbox{mainbox}[1]{
  colframe=maincolor,
  base={#1}
}

% Sub concept box (gray accent)
% Use for: examples, remarks, minor concepts
\newtcolorbox{subbox}[1]{
  colframe=black!20!white,
  base={#1}
}

% Alternative colors (uncomment to use)
% \definecolor{maincolor}{rgb}{0.2, 0.4, 0.8}   % Blue
% \definecolor{maincolor}{rgb}{0.8, 0.2, 0.2}   % Red
% \definecolor{maincolor}{rgb}{0.5, 0.2, 0.6}   % Purple

% ============================================================================
% MATHEMATICS
% ============================================================================
\usepackage{amsmath, amsthm, mathtools, amssymb}
\usepackage{amsfonts, amstext}
\allowdisplaybreaks  % Allow page breaks in equations

% Common math shortcuts (customize as needed)
\def\limn{\lim_{n\to \infty}}
\def\limxo{\lim_{x\to 0}}
\def\limxi{\lim_{x\to\infty}}
\def\limxn{\lim_{x\to-\infty}}
\def\sumk{\sum_{k=1}^\infty}
\def\sumn{\sum_{n=0}^\infty}

% Number sets
\def\N{\mathbb{N}}  % Natural numbers
\def\Z{\mathbb{Z}}  % Integers
\def\Q{\mathbb{Q}}  % Rationals
\def\R{\mathbb{R}}  % Reals
\def\C{\mathbb{C}}  % Complex

% Calculus
\def\dx{\,dx}
\def\dt{\,dt}
\def\dy{\,dy}
\newcommand{\ddx}[1]{\frac{d#1}{dx}}
\newcommand{\ddn}[2]{\frac{d#1}{d#2}}
\newcommand{\pdn}[2]{\frac{\partial #1}{\partial #2}}

% Big O notation
\def\oof{\mathcal{O}}

% ============================================================================
% TABLES
% ============================================================================
\usepackage{tabularx, multirow, booktabs}
\renewcommand*{\arraystretch}{1.5}  % More spacing in tables

% ============================================================================
% LISTS
% ============================================================================
\usepackage{enumitem}
\setitemize{itemsep=0.5pt, topsep=2pt}
\setenumerate{itemsep=0.75pt, topsep=2pt}

% ============================================================================
% GRAPHICS & TIKZ
% ============================================================================
\usepackage{graphicx}
\usepackage{tikz}
\usetikzlibrary{
  positioning,
  calc,
  fit,
  matrix,
  decorations,
  angles,
  quotes,
  babel,
  3d,
  arrows.meta,
  shadings,
  shapes.geometric
}

\tikzset{
  >=stealth,
  % Add custom TikZ styles here
}

% ============================================================================
% HYPERLINKS
% ============================================================================
\usepackage{hyperref}
\hypersetup{
  colorlinks=true,
  linkcolor=maincolor,
  urlcolor=maincolor
}

% ============================================================================
% ADDITIONAL PACKAGES
% ============================================================================
\usepackage{centernot}   % For \centernot command
\usepackage{marvosym}    % Additional symbols
\usepackage{wasysym}     % Additional symbols

% ============================================================================
% DOCUMENT METADATA
% ============================================================================
\title{Cheatsheet ANA-F}
\author{Thomas Schnyder}
\date{September 2026}

% Table of contents depth
\setcounter{tocdepth}{4}
\setcounter{secnumdepth}{4}

% ============================================================================
% DOCUMENT START
% ============================================================================
\begin{document}

\maketitle

% ============================================================================
% SECTION 1: TRIGONOMETRISCHE FUNKTIONEN
% ============================================================================
\section{Trigonometrische Funktionen}

\begin{mainbox}{Definition am Einheitskreis}
Punkt $P$ auf dem Einheitskreis mit Winkel $x$ (zur pos. $x$-Achse):
$$P = (\cos x,\ \sin x)$$
$$\tan x = \frac{\sin x}{\cos x}, \qquad \cot x = \frac{\cos x}{\sin x} = \frac{1}{\tan x}$$
\textbf{Bogenmass:} $x = \dfrac{\pi}{180^\circ} \cdot \alpha$ \quad ($360^\circ = 2\pi$)
\end{mainbox}

\begin{mainbox}{Wichtige Werte}
\renewcommand*{\arraystretch}{1.6}
\setlength{\tabcolsep}{4pt}
\begin{center}
\begin{tabular}{c|ccccc}
$x$ & $0$ & $\frac{\pi}{6}$ & $\frac{\pi}{4}$ & $\frac{\pi}{3}$ & $\frac{\pi}{2}$ \\
\hline
$\sin x$ & $0$ & $\frac{1}{2}$ & $\frac{\sqrt2}{2}$ & $\frac{\sqrt3}{2}$ & $1$ \\
$\cos x$ & $1$ & $\frac{\sqrt3}{2}$ & $\frac{\sqrt2}{2}$ & $\frac{1}{2}$ & $0$ \\
$\tan x$ & $0$ & $\frac{\sqrt3}{3}$ & $1$ & $\sqrt3$ & -- \\
\end{tabular}
\end{center}
\textbf{Merkregel:} $\sin x = \frac{\sqrt{n}}{2}$ für $n = 0,1,2,3,4$; $\cos$ rückwärts.
\end{mainbox}

\begin{mainbox}{Eigenschaften}
\renewcommand*{\arraystretch}{1.5}
\setlength{\tabcolsep}{3pt}
\begin{tabular}{l|cccc}
 & $\sin$ & $\cos$ & $\tan$ & $\cot$ \\
\hline
$W$ & $[-1,1]$ & $[-1,1]$ & $\R$ & $\R$ \\
Periode & $2\pi$ & $2\pi$ & $\pi$ & $\pi$ \\
Symm. & ung. & ger. & ung. & ung. \\
Nullst. & $k\pi$ & $\frac{\pi}{2} + k\pi$ & $k\pi$ & $\frac{\pi}{2} + k\pi$ \\
\end{tabular}
$D_{\sin} = D_{\cos} = \R$, \ $D_{\tan} = \R \setminus \{\frac{\pi}{2} + k\pi\}$,
$D_{\cot} = \R \setminus \{k\pi\}$ ($k \in \Z$). Polstellen von $\tan$/$\cot$ = Nullstellen von $\cos$/$\sin$.
\end{mainbox}

\begin{mainbox}{Wichtige Formeln}
\textbf{Trigonometrischer Pythagoras:}
$$\boxed{\sin^2 x + \cos^2 x = 1}$$
\textbf{Additionstheoreme:}
\begin{align*}
\sin(x \pm y) &= \sin x \cos y \pm \cos x \sin y \\
\cos(x \pm y) &= \cos x \cos y \mp \sin x \sin y \\
\tan(x \pm y) &= \frac{\tan x \pm \tan y}{1 \mp \tan x \tan y}
\end{align*}
\textbf{Doppelwinkel} ($y = x$):
\begin{align*}
\sin(2x) &= 2 \sin x \cos x \\
\cos(2x) &= \cos^2 x - \sin^2 x \\
         &= 1 - 2\sin^2 x = 2\cos^2 x - 1
\end{align*}
\end{mainbox}

\begin{subbox}{Symmetrien \& Verschiebungen}
\begin{align*}
\sin(-x) &= -\sin x & \cos(-x) &= \cos x \\
\sin(\pi - x) &= \sin x & \cos(\pi - x) &= -\cos x
\end{align*}
\begin{align*}
\sin\!\left(x + \tfrac{\pi}{2}\right) &= \cos x \\
\cos\!\left(x + \tfrac{\pi}{2}\right) &= -\sin x
\end{align*}
\end{subbox}

\begin{subbox}{Allgemeine Sinusfunktion}
$$y = A \cdot \sin(\omega x + \varphi)$$
Amplitude $|A|$, Periode $p = \frac{2\pi}{\omega}$, Verschiebung um
$-\frac{\varphi}{\omega}$ (nach links, falls $\varphi > 0$).
\end{subbox}

\begin{mainbox}{Ableitungen}
\begin{align*}
(\sin x)' &= \cos x \\
(\cos x)' &= -\sin x \\
(\tan x)' &= \frac{1}{\cos^2 x} = 1 + \tan^2 x \\
(\cot x)' &= -\frac{1}{\sin^2 x} = -(1 + \cot^2 x)
\end{align*}
$\tan'$ via Quotientenregel: $\frac{\cos x \cos x - \sin x (-\sin x)}{\cos^2 x}
= \frac{1}{\cos^2 x}$.
\end{mainbox}

\subsection{Arkusfunktionen}

\begin{mainbox}{Umkehrfunktionen (Hauptwerte)}
Nur auf einem \textbf{Monotonieintervall} umkehrbar $\Rightarrow$ eingeschränkter
Wertebereich:
\renewcommand*{\arraystretch}{1.5}
\setlength{\tabcolsep}{4pt}
\begin{center}
\begin{tabular}{l|c|c}
 & $D$ & $W$ \\
\hline
$\arcsin x$ & $[-1, 1]$ & $[-\frac{\pi}{2}, \frac{\pi}{2}]$ \\
$\arccos x$ & $[-1, 1]$ & $[0, \pi]$ \\
$\arctan x$ & $\R$ & $(-\frac{\pi}{2}, \frac{\pi}{2})$ \\
$\operatorname{arccot} x$ & $\R$ & $(0, \pi)$ \\
\end{tabular}
\end{center}
$y = \arcsin x \iff x = \sin y$ (mit $y \in W$).\\
$\arcsin x + \arccos x = \frac{\pi}{2}$, \ $\arctan x + \operatorname{arccot} x = \frac{\pi}{2}$
\end{mainbox}

\begin{mainbox}{Ableitungen der Arkusfunktionen}
\begin{align*}
(\arcsin x)' &= \frac{1}{\sqrt{1 - x^2}} \\
(\arccos x)' &= -\frac{1}{\sqrt{1 - x^2}} \\
(\arctan x)' &= \frac{1}{1 + x^2} \\
(\operatorname{arccot} x)' &= -\frac{1}{1 + x^2}
\end{align*}
\end{mainbox}

\begin{subbox}{Denkprozess: Herleitung $(\arcsin x)'$}
\begin{enumerate}
    \item $y = \arcsin x \iff \sin y = x$
    \item Implizit ableiten: $\cos y \cdot y' = 1 \Rightarrow y' = \frac{1}{\cos y}$
    \item $\cos y = \sqrt{1 - \sin^2 y}$ (da $\cos y \geq 0$ auf $[-\frac{\pi}{2}, \frac{\pi}{2}]$)
    \item $\Rightarrow y' = \frac{1}{\sqrt{1 - x^2}}$
\end{enumerate}
Analog $\arctan$: $\tan y = x \Rightarrow (1 + \tan^2 y)\,y' = 1 \Rightarrow y' = \frac{1}{1 + x^2}$.
\end{subbox}

\begin{subbox}{Beispiel mit Kettenregel}
$$\left(\arctan(x^2)\right)' = \frac{1}{1 + x^4} \cdot 2x = \frac{2x}{1 + x^4}$$
\end{subbox}

\subsection{Trigonometrische Gleichungen}

\begin{mainbox}{Alle Lösungen ($k \in \Z$)}
\renewcommand*{\arraystretch}{1.5}
\begin{itemize}
    \item $\sin z = a$: \ $z = \arcsin a + 2k\pi$ \ oder \ $z = \pi - \arcsin a + 2k\pi$
    \item $\cos z = a$: \ $z = \pm \arccos a + 2k\pi$
    \item $\tan z = a$: \ $z = \arctan a + k\pi$
\end{itemize}
(für $\sin$/$\cos$ nur lösbar, falls $|a| \leq 1$)
\end{mainbox}

\begin{subbox}{Vorgehen}
\begin{enumerate}
    \item Mit Formeln (Pythagoras, Doppelwinkel) auf \emph{eine} Funktion bringen.
    \item Ggf. substituieren ($u = \sin x$) und algebraisch lösen.
    \item Grundlösung(en) mit Arkusfunktion, dann \emph{alle} Lösungen über
          Symmetrie + Periode.
    \item Zurücksubstituieren, nach $x$ auflösen, ggf. auf Intervall einschränken.
\end{enumerate}
\end{subbox}

\begin{subbox}{Beispiel: $\sin(2x) = \frac{1}{2}$}
$2x = \frac{\pi}{6} + 2k\pi$ \ oder \ $2x = \frac{5\pi}{6} + 2k\pi$\\
$\Rightarrow x = \frac{\pi}{12} + k\pi$ \ oder \ $x = \frac{5\pi}{12} + k\pi$
\end{subbox}

\begin{subbox}{Beispiel: $2\sin^2 x - \sin x - 1 = 0$}
$u = \sin x$: $2u^2 - u - 1 = 0 \Rightarrow u_1 = 1$, $u_2 = -\frac{1}{2}$
\begin{itemize}
    \item $\sin x = 1 \Rightarrow x = \frac{\pi}{2} + 2k\pi$
    \item $\sin x = -\frac{1}{2} \Rightarrow x = -\frac{\pi}{6} + 2k\pi$ \ oder \
          $x = \frac{7\pi}{6} + 2k\pi$
\end{itemize}
\end{subbox}

\begin{subbox}{Häufige Fehler \& Tipps}
\begin{itemize}
    \item Taschenrechner auf \textbf{RAD} stellen.
    \item Arkusfunktion liefert nur \emph{eine} Lösung -- zweite Lösung
          (Symmetrie) und Periode nicht vergessen.
    \item Nicht durch $\sin x$ oder $\cos x$ dividieren (Lösungen gehen verloren)
          -- stattdessen ausklammern.
    \item $\sin^2 x = (\sin x)^2 \neq \sin(x^2)$.
    \item $\arcsin x \neq \frac{1}{\sin x}$ (Umkehrfunktion, nicht Kehrwert).
    \item Periode von $\sin(\omega x)$ ist $\frac{2\pi}{\omega}$ $\Rightarrow$
          bei $\sin(2x)$: $+k\pi$.
\end{itemize}
\end{subbox}

% ============================================================================
% TIPS FOR AVOIDING OVERFULL BOXES
% ============================================================================
%
% 1. Keep formulas concise - split long equations over multiple lines
% 2. Use shorter variable names in inline math
% 3. Abbreviate table headers
% 4. Use align* environment for multi-line equations
% 5. Test compile frequently to catch overflow early
%
% Example of splitting a long formula:
% Instead of: $$a + b + c + d + e = f + g + h + i + j$$
% Use:
% \begin{align*}
% a + b + c + d + e &= f + g + h \\
%                   &\quad + i + j
% \end{align*}
%
% ============================================================================

\end{document}